\subsection{About the Course}
\label{sec:about-the-course}
\slidesrc{Intro: 1--3; Terminology introduction: 2}

The course (taught by Viviana Mascardi, University of Genoa) loosely follows the textbook
\emph{Speech and Language Processing} (3rd edition draft) by Dan Jurafsky and James H.
Martin, freely available at \url{https://web.stanford.edu/~jurafsky/slp3/}. Some slides
integrate those associated with the book. There is no single reference book, however:
the slides, plus the pointers to material on the web and on the course's AulaWeb page,
are enough to understand the subject, even without attending the lessons.

\begin{itemize}[nosep]
  \item A large share of the lessons covers \textbf{``Annotating Linguistic Structure''}
        (Volume~III of the book), extended with lessons on the birth of human language
        (Section~\ref{sec:birth-of-language}), semantics and pragmatics.
  \item The remaining lessons cover \textbf{Large Language Models} (Volume~I).
  \item Topics of Volume~I already covered in other Machine Learning courses (logistic
        regression, neural networks, RNNs, \ldots) are \emph{not} repeated.
\end{itemize}

\begin{nbfact}{Exam}
The exam has two parts: a \textbf{written exam} and an \textbf{individual project}.
Students can propose the project themselves, so they are free to experiment in a
practical way with the topics they liked most.
\end{nbfact}

The (preliminary) programme for 2026--2027 is, in order: introduction and regular
expressions; syntax at the word level (segmentation, stop words, TF-IDF, stemming,
lemmatisation); minimum edit distance, the sentence level and rule-based POS tagging;
grammars and statistical (Hidden Markov Model) POS tagging; semantics (distributional,
frame, model-theoretic and lexical semantics, a comparison among types of semantics,
ontologies and the Semantic Web); practical labs on NLP libraries, text summarisation,
chatbots and agentic AI; pragmatics; a wrap-up on classical NLP; and finally BERT, GPT
and LLMs. Section~\ref{sec:course-map} shows how these topics fit together.
